\section{Revestimentos Avançados e Manufatura Funcional}\label{sec:revestimentos_manufatura_funcional}

A engenharia de materiais para solicitações térmicas extremas transitou de uma abordagem restrita a ligas maciças monolíticas para o desenvolvimento integrado de Sistemas de Proteção Térmica (TPS). Nesse paradigma híbrido, ligas refratárias de alta entropia operam como substratos estruturais primários, enquanto revestimentos superficiais funcionais fornecem a barreira indispensável contra o ataque oxidativo e a degradação aerotérmica.

\subsection{Revestimentos e Hibridização Superficial}\label{subsec:revestimentos_hibridizacao}

Superfícies desprotegidas de ligas convencionais sofrem oxidação descontrolada acima de $1100~^\circ\mathrm{C}$. Em contrapartida, as RHEAs projetadas para barreiras térmicas (TBCs) promovem a nucleação de camadas contínuas de óxido de alumínio ($\alpha\text{-Al}_2\text{O}_3$), isolando termicamente o substrato e reduzindo sensivelmente a taxa de ganho de massa. Para suportar regimes ainda mais severos, em temperaturas próximas a $2100~^\circ\mathrm{C}$, revestimentos de silicietos dopados com elementos reativos (como o ítrio) introduzem capacidade autorregenerativa (\emph{self-healing}). Sob o fluxo aerotérmico, esses compostos formam uma película vítrea de borossilicato de alta viscosidade que flui sobre a superfície e sela autonomamente microfissuras e defeitos superficiais, bloqueando a difusão do oxigênio gasoso antes que ocorra a degradação estrutural do núcleo refratário.

\subsection{Compósitos Multifásicos e Reforço Estrutural}\label{subsec:compositos_multifasicos}

A superação dos limites mecânicos intrínsecos aos metais puros impulsionou o desenvolvimento de compósitos multifásicos que aliam a tenacidade da matriz metálica refratária à rigidez e estabilidade de fases cerâmicas. A introdução e dispersão homogênea de nanopartículas de carbetos e nitretos de ultra-alta temperatura (UHTCs) estabelece obstáculos densos à movimentação de discordâncias e à migração de contornos de grão. Essa microestrutura complexa eleva substancialmente a resistência à fluência em temperaturas superiores a $1500~^\circ\mathrm{C}$ e mitiga a fragilidade excessiva à temperatura ambiente por meio da modulação de tensões locais nas interfaces intermetálicas.

\subsection{Rotas de Processamento e Metalurgia do Pó}\label{subsec:manufatura_avancada}

A fundição tradicional por arco elétrico comumente induz macrossegregação química, porosidades e tensões residuais severas em ligas refratárias multicomponentes. Por essa razão, a literatura converge amplamente para técnicas avançadas de metalurgia do pó. A Sinterização por Plasma de Centelha (\emph{Spark Plasma Sintering} -- SPS) destaca-se ao empregar pulsos de corrente elétrica combinados a pressão uniaxial, consolidando os pós metálicos em ciclos de poucos minutos. Essa taxa de aquecimento ultrarrápida suprime o crescimento excessivo de grão e preserva os benefícios do endurecimento refinado por Hall-Petch. Complementarmente, a Prensagem Isostática a Quente (\emph{Hot Isostatic Pressing} -- HIP) aplica calor e pressão hidrostática simultânea em atmosfera inerte para colapsar microporosidades internas e homogeneizar quimicamente a liga. Para sistemas que exigem elevada ductilidade, processos subsequentes de deformação plástica severa, como a laminação termomecânica a quente, refinam subgrãos e quebram fases frágeis, melhorando a tolerância ao dano mecânico à temperatura ambiente.

\subsection{Manufatura Aditiva e Simulação Acelerada por Inteligência Artificial}\label{subsec:impresao_3d_ia}

A viabilização de geometrias aeroespaciais avançadas, como bordos de ataque e bordos de fuga com canais internos de resfriamento regenerativo, consolidou a Manufatura Aditiva por Fusão em Leito de Pó a Laser (\emph{Laser Powder Bed Fusion} -- LPBF) como vetor estratégico. As altíssimas taxas de resfriamento inerentes ao processo a laser ($10^4$ a $10^6~\mathrm{K/s}$) suprimem a segregação elementar e geram microestruturas celulares ultrafinas, conferindo elevada resistência mecânica diretamente após a conformação. Simultaneamente, a velocidade de exploração de novos espaços de composição foi radicalmente transformada pela integração entre modelagem termodinâmica via CALPHAD e algoritmos de aprendizado de máquina (\emph{Machine Learning}). Em vez de campanhas experimentais baseadas em tentativa e erro, modelos preditivos rastreiam milhões de estequiometrias virtuais para projetar ligas termodinamicamente estáveis e com baixa propensão a fases intermetálicas deletérias, acelerando de forma sem precedentes o ciclo de maturidade tecnológica desses materiais para voos hipersônicos.
